\documentclass[12pt]{article}

\usepackage[T1]{fontenc}
\usepackage[utf8]{inputenc}
\usepackage[brazil]{babel}
\usepackage{lmodern}
\usepackage{microtype}
\usepackage[a4paper,margin=2.5cm]{geometry}
\usepackage{setspace}
\usepackage{indentfirst}

\setstretch{1.15}

\title{Métodos e Problemas na Teoria dos Ciclos Econômicos \thanks{Texto Traduzido livremente por Luiz Mario Andrade com o auxílio da ferramenta Chat GPT 5.4}}
\author{Robert E. Lucas, Jr.}
\date{Novembro de 1980}

\begin{document}

\maketitle

\let\thefootnote\relax\footnotetext{* Este artigo foi preparado para o Seminário do American Enterprise Institute sobre Expectativas Racionais, realizado em 1 de fevereiro de 1980, em Washington, D.C. Muitas das ideias nele contidas foram desenvolvidas sob o estímulo do curso de Don Patinkin sobre a História do Pensamento Monetário, ministrado na Universidade de Chicago, no inverno de 1979. Allan Drazen, Sherwin Rosen e Nasser Saidi forneceram críticas muito úteis a uma versão preliminar. Uma versão deste artigo intitulada "Política Econômica e o Ciclo Econômico" foi apresentada na Universidade Estadual de Ohio como a Palestra sobre Moeda, Crédito e Bancos, em 8 de maio de 1980.
\par ROBERT E. LUCAS, JR. é professor de economia na Universidade de Chicago.
\par \copyright 1980 American Enterprise Institute for Public Policy Research
\par JOURNAL OF MONEY, CREDIT, AND BANKING, vol. 12, no. 4 (Novembro 1980, Parte 2)}

\section{INTRODUÇÃO}

UMA DAS FUNÇÕES da economia teórica é fornecer sistemas econômicos artificiais, plenamente articulados, que possam servir como laboratórios nos quais políticas que seriam proibitivamente caras para experimentar em economias reais possam ser testadas a um custo muito menor. Para desempenhar bem essa função, é essencial que a economia do "modelo" artificial seja distinguida da forma mais nítida possível, na discussão, das economias reais. Na medida em que houver confusão entre declarações de opinião sobre como acreditamos que as economias reais reagiriam a políticas específicas e declarações de fatos verificáveis sobre como o modelo reagirá, a teoria não estará sendo usada de forma eficaz para nos ajudar a ver quais opiniões sobre o comportamento das economias reais são precisas e quais não são. É nesse sentido que a insistência na "realidade" de um modelo econômico subverte sua utilidade potencial para pensar sobre a realidade. Qualquer modelo que seja suficientemente articulado para dar respostas claras às perguntas que lhe fazemos será necessariamente artificial, abstrato, patentemente "irreal".

Ao mesmo tempo, nem todos os modelos bem articulados serão igualmente úteis. Embora estejamos interessados em modelos porque acreditamos que eles podem nos ajudar a entender questões sobre as quais somos atualmente ignorantes, precisamos testá-los como imitações úteis da realidade, sujeitando-os a choques para os quais temos relativa certeza de como as economias reais, ou partes delas, reagiriam. Quanto mais dimensões nas quais o modelo imita as respostas que as economias reais dão a perguntas simples, mais confiamos em suas respostas a perguntas mais difíceis. É nesse sentido que mais "realismo" em um modelo é claramente preferível a menos.

Nesta visão geral da natureza da teoria econômica, portanto, uma "teoria" não é uma coleção de afirmações sobre o comportamento da economia real, mas sim um conjunto explícito de instruções para construir um sistema paralelo ou análogo — uma economia mecânica de imitação. Um "bom" modelo, desse ponto de vista, não será exatamente mais "real" do que um pobre, mas fornecerá melhores imitações. É claro que o que se entende por uma "melhor imitação" dependerá das perguntas específicas para as quais se deseja respostas.\footnote{Não conheço os antecedentes desta visão da teoria como análogo físico, nem tenho uma ideia clara de quão amplamente compartilhada ela é entre os economistas. Um ancestral imediato da minha afirmação condensada é [43].}

Neste artigo, desejo revisar alguns desenvolvimentos recentes na teoria dos ciclos econômicos, adotando o ponto de vista sugerido acima. Sob essa ótica, seria de se esperar que os desenvolvimentos surgissem de dois tipos bastante diferentes de forças externas às subdisciplinas da economia monetária ou da teoria dos ciclos econômicos. Dessas forças, a mais importante, acredito, nesta área e na economia em geral, consiste em desenvolvimentos puramente técnicos que aumentam nossas capacidades de construir economias análogas. Aqui eu incluiria tanto melhorias em métodos matemáticos quanto melhorias na capacidade computacional. A negligência na história tradicional da doutrina dessa força de mudança em nosso pensamento é uma omissão séria e contribui para o sentimento comum, mas equivocado, de que tudo já foi dito antes ou de que "está tudo em Marshall". O mundo de Marshall era suficientemente parecido com o nosso e Marshall era um observador astuto o suficiente de seu mundo para que seja difícil fazer observações gerais sobre nossa economia que não tenham um precedente próximo em alguma observação marshalliana sobre a dele. Nossa capacidade de construir economias análogas é, no entanto, muito maior, de modo que temos a capacidade de estudar em detalhes as interações de mercado sobre as quais Marshall só podia conjeturar.

A segunda fonte de desenvolvimentos teóricos são as mudanças nas perguntas que queremos que os modelos respondam, ou nos fenômenos que desejamos entender ou explicar. Para o jornalista, cada ano traz novos fenômenos sem precedentes, exigindo novas teorias sem precedentes (onde "teoria" equivale a uma descrição dos novos fenômenos juntamente com a afirmação de que eles são novos). Como há um sentido óbvio em que essa visão carrega alguma verdade, não tentarei refutá-la. Argumentei em outro lugar [23] que é de nosso interesse adotar exatamente o ponto de vista oposto no estudo dos ciclos econômicos, ou o mais próximo de um ponto de vista oposto que pudermos, e manterei essa atitude a seguir. A Grande Depressão, no entanto, permanece uma barreira formidável a uma aplicação completamente inflexível da visão de que os ciclos econômicos são todos iguais.

Existe, é claro, uma terceira fonte de desenvolvimentos em nossa compreensão dos ciclos econômicos: as atividades de economistas especializados no campo. Fazer as conexões entre as inovações técnicas descobertas por nossos colegas e as perguntas difíceis lançadas a nós pelo mundo real não é tarefa fácil, e não desejo minimizar a importância desses esforços. No entanto, penso que tanto historiadores amadores quanto profissionais tenderam a ir longe demais ao tentar entender os desenvolvimentos na economia monetária em termos inteiramente internos à subdisciplina, e que alguma inclinação na outra direção pode, portanto, ser útil.

Na próxima seção do artigo, revisarei o que me parecem ser as principais características da Revolução Keynesiana a partir da visão apresentada acima. Há, eu sei, um sentimento crescente de que tal "sacudir de esqueletos" está se tornando cansativo e antiquado. No entanto, o avanço de constructos teóricos dessa era como se fossem fatos a serem explicados continua a ser um modo padrão de debate em macroeconomia, sugerindo a existência de ossos ainda mais profundamente enterrados. O restante do artigo será uma tentativa de diagnosticar desenvolvimentos mais recentes e talvez até extrapolar um pouco para o futuro.

\section{TEORIA DOS CICLOS ECONÔMICOS ATÉ KEYNES}

A teoria dos ciclos econômicos (como distinta da economia monetária) é principalmente um produto do século XX. Na maior parte, os principais economistas do século XIX deixaram as flutuações de curto prazo de lado para focar a atenção em outras questões. A ideia geral subjacente (poder-se-ia chamá-la de uma "hipótese de taxa natural") deve ter sido que se poderia entender os principais determinantes dos níveis médios e taxas de crescimento da atividade econômica sem entender as flutuações.

A tarefa de tornar preciso exatamente quais aspectos da vida econômica haviam sido deixados de lado por essa estratégia frutífera do século XIX foi assumida no início do século atual por Wesley Mitchell. Mitchell buscou uma definição empírica de ciclos econômicos por meio da exclusão sistemática daqueles movimentos nas séries temporais econômicas que pareciam passíveis de explicação pela teoria então existente: o nível geral e o padrão de crescimento na atividade econômica, e os movimentos em séries individuais que pareciam surgir de condições de oferta ou demanda específicas para mercados individuais. Nenhum desses princípios de exclusão é isento de ambiguidade, por isso Mitchell procedeu com cautela e experimentalmente, testando uma ampla variedade de técnicas de sumarização de dados na mais ampla coleção disponível de séries.

Por duas razões, é fácil esquecer o caráter notável das regularidades que Mitchell conseguiu descobrir e documentar. Por um lado, vivemos com elas por tanto tempo que parecem não tanto o produto de uma forma imaginativa e abstrata de organizar as séries temporais econômicas, mas simplesmente "fatos" que "todo mundo conhece". Por outro lado, são regularidades que, do ponto de vista mais amplamente adotado desde a década de 1930, não são especialmente dignas de nota ou sugestivas. A descoberta central, é claro, foi a semelhança de todos os ciclos de paz uns com os outros, uma vez controlada a variação na duração, no sentido de que cada ciclo exibe aproximadamente o mesmo padrão de comovimentos entre as variáveis que os outros.

Não surpreendentemente, essas descobertas empíricas (guiadas teoricamente no sentido que indiquei), corroboradas por outras evidências e uma variedade de impressões menos sistematicamente organizadas, estimularam uma grande quantidade de trabalho teórico. A ideia de que se poderia, com uma base empírica firme, falar de algo como um "ciclo econômico típico", dividido em estágios invariantes em caráter (se não em duração), sugeria que uma parte substancial das flutuações observadas poderia ser explicável em um nível bastante abstrato, ou "simples", com uma única explicação teórica do (ou seja, de todos os) ciclo(s) econômico(s). Este é o princípio organizador que Gottfried Haberler empregou, por exemplo, em \textit{Prosperity and Depression} [11]. Que uma variedade tão ampla de teorias, como as discutidas naquele livro, pudesse ser ajustada com tanta facilidade em um padrão lógico tão apertado é um testemunho, creio eu, do poder da abstração na ideia de um ciclo econômico único ou típico.

Outro exemplo de teoria dos ciclos econômicos pré-keynesiana (aqui suponho que estou abusando imperdoavelmente da linguagem), que é útil para dois dos meus propósitos, é o \textit{Treatise on Money} [16] de John Maynard Keynes. A comparação deste trabalho com a \textit{General Theory} [17] é útil tanto para ilustrar a maneira como os limites em nossa capacidade técnica de construir teoria explícita limitam nossa capacidade de pensar produtivamente sobre os fenômenos, quanto para ilustrar a extensão em que o pensamento na economia monetária está sujeito a choques externos ou do "mundo real".

O \textit{Treatise} é por vezes descrito (e o foi até pelo próprio Keynes\footnote{No prefácio da \textit{General Theory}, pp. vi-vii.}) como uma espécie de tateio inicial e desajeitado em direção às ideias apresentadas na \textit{General Theory}. Suponho que deva haver alguma verdade psicológica nessa visão, mas enfatizar esse aspecto do \textit{Treatise} leva a uma leitura muito forçada do que me parece ser um exemplo bastante direto do pensamento pré-Depressão sobre ciclos econômicos. O objetivo principal do livro é tentar entender as flutuações na atividade econômica em torno de uma tendência secular na qual as magnitudes reais são determinadas pelas considerações reais da teoria do valor neoclássica e na qual os preços nominais são governados pela teoria quantitativa da moeda. Keynes estava convencido, corretamente acredito, de que tentar discutir flutuações em termos de flutuações na velocidade não seria produtivo e buscou, em vez disso, um ponto de vista que enfatizasse mudanças na composição dos gastos ao longo do ciclo. Ele então mostrou que esse ponto de vista poderia ser reconciliado com uma visão quantitativista de mudanças de longo prazo.

O sistema contábil, ou conjunto de convenções notacionais, embutido em suas "equações fundamentais", foi projetado para facilitar essa reconciliação e serviu bem a essa função. Além disso, no entanto, o aparato de Keynes não podia ir. Ele afirma uma identidade e a discute; então move um termo da esquerda para a direita e discute o resultado dessa operação; define uma nova variável em termos de outras previamente definidas e discute isso; então fala sobre a identidade reafirmada em termos dessa nova variável; e assim por diante. O problema não é que as ideias subjacentes sejam triviais, embora a álgebra certamente o seja. Pelo contrário, o livro lida de maneira inteligente com os problemas fundamentais que os ciclos econômicos levantam. A dificuldade é que Keynes não possui aparato para lidar com esses problemas. Embora os discuta verbalmente tão bem quanto seus contemporâneos, nem ele nem ninguém estava tecnicamente equipado o suficiente para elevar a discussão a um nível mais nítido ou mais produtivo.

O início da Grande Depressão não contribuiu em nada para melhorar o equipamento de Keynes para entender o ciclo econômico, visto como uma sequência recorrente de booms e depressões. Em vez disso, permitiu-lhe reformular o próprio problema como sendo o de explicar o nível de produto e emprego em um ponto no tempo, em oposição a explicar um padrão particular repetido nas séries temporais. Assim reformulado, o problema poderia ser estudado produtivamente simplesmente descartando uma equação da teoria do equilíbrio estático (a curva de oferta de trabalho) em contraste com a tarefa muito mais difícil, empreendida no \textit{Treatise}, de suplementar esta teoria estática com uma dinâmica de curto prazo adequada. Este problema mais simples era um no qual o progresso poderia ser feito no nível de análise marshalliano no qual Keynes era um mestre.\footnote{Novamente, do prefácio da \textit{General Theory}, p. viii: "A dificuldade reside não nas novas ideias, mas em escapar das antigas." Mas seria necessário mais trabalho, eu sei, do que localizar uma única citação para convencer um cético sobre este ponto.}

Ao ler a \textit{General Theory} desta maneira, estou naturalmente seguindo a exegese clássica de John Hicks [14] (bem como a visão inicial de Hicks da \textit{General Theory} como "economia de crise") e o passo pioneiro de Franco Modigliani [29] em direção a uma "síntese neoclássica". Há, certamente, muito de interesse na \textit{General Theory} que não é capturado nem no diagrama de Hicks nem no sistema de equações de Modigliani, um fato que levou Axel Leijonhufvud (e outros, talvez até os próprios Hicks e Modigliani) a ver a "economia keynesiana", que foi posteriormente baseada principalmente nessas interpretações iniciais, como uma espécie de vulgarização da \textit{General Theory}. Embora haja alguma verdade, vigorosamente desenvolvida na monografia de Leijonhufvud [18], nesta visão, ela perde o que acredito ser a verdade mais essencial, enfatizada em minha introdução, de que o progresso no pensamento econômico significa obter modelos econômicos análogos, abstratos, cada vez melhores, e não melhores observações verbais sobre o mundo.

A \textit{General Theory} é, com certeza, uma mina de observações agudas e bem formuladas sobre as dificuldades de conduzir os próprios negócios em um mundo incerto. Talvez existam algumas, embora eu não gostasse de ter a tarefa de documentar isso, que sejam centrais para o comportamento do ciclo econômico e não prefiguradas em, digamos, Mitchell [28], ou mesmo um século antes nos escritos de Henry Thornton. Certamente o provável papel central das expectativas de lucro e do comportamento do investimento foi enfatizado por virtualmente todos os economistas que dedicaram mais do que um pensamento superficial aos ciclos econômicos. Economistas que consideram o estilo de Keynes atraente continuarão a usar seus escritos como Dennis Robertson usou os de Lewis Carroll, mas certamente há mais na natureza cumulativa da economia do que isso!

Extrair da \textit{General Theory} um método gráfico simples para pensar sobre a determinação da renda nacional não é, acredito, vulgarizar sua contribuição. A vulgaridade na economia seria mais apropriadamente definida como criticar ou caricaturar um modelo abstrato (e, portanto, potencialmente útil) porque ele omite algo.

\section{A SÍNTESE NEOCLÁSSICA}

A \textit{General Theory} foi, se o argumento da seção anterior estiver correto, mais bem-sucedida que o \textit{Treatise} não por causa de avanços teóricos dentro da economia monetária, mas porque o evento da Grande Depressão permitiu que Keynes reformulasse o problema colocado pelo ciclo econômico em uma forma tal que os métodos teóricos à sua disposição admitissem progresso genuíno. Foi um acidente histórico afortunado que, aproximadamente ao mesmo tempo e por razões não relacionadas aos eventos econômicos contemporâneos, ocorreram avanços técnicos na teoria estatística e econômica, os quais transformaram a "economia keynesiana" em algo muito diferente de, e muito mais frutífero do que, qualquer coisa que o próprio Keynes tivesse previsto.

Um desses avanços, que se pode datar do trabalho de Tinbergen [47] ou talvez de Slutsky [44], foi a ideia de que se poderia descrever uma economia como um sistema de equações de diferenças perturbadas estocasticamente, cujos parâmetros poderiam ser estimados a partir de séries temporais reais. De fato, ao descrever uma "teoria do ciclo econômico" como uma "economia de imitação plenamente articulada", pressupus que o objetivo desses primeiros econometricistas se tornou agora propriedade comum. O rápido progresso dos modelos econométricos em direção a algo que parecia próximo da perfeição imitativa foi, como argumentei em outro lugar [22], ilusório. Sargent e eu [26] argumentamos posteriormente em detalhes que sistemas análogos úteis não serão encontrados entre variações dentro da classe desses modelos econométricos pioneiros. As questões abordadas nesses artigos anteriores ainda me parecem estar no centro da questão do que podemos esperar de uma teoria dos ciclos econômicos, mas não há razão para retomá-las aqui.

O segundo desses avanços, como o primeiro, tem raízes cujo deslindamento desafiaria a carreira de um estudioso objetivo. Em vez disso, seguirei um curso subjetivo, identificando esse desenvolvimento com o livro \textit{Foundations of Economic Analysis} [36] de Paul Samuelson e aconselhando o historiador mais sério a buscar as copiosas referências naquele texto. Ao fazê-lo, sigo a prática de Don Patinkin em seu \textit{Money, Interest, and Prices} [32], talvez a versão mais refinada e influente do que quero dizer com o termo "síntese neoclássica". Samuelson avançou, em primeiro lugar, os principais ingredientes para uma teoria matematicamente explícita do equilíbrio geral: um sistema artificial no qual famílias e empresas resolvem conjuntamente problemas de maximização "estáticos" explícitos, tomando os preços como dados parametricamente. Ele deu como certo que tais equilíbrios não eram vazios no sentido matemático, uma suposição que foi confirmada posteriormente sob condições bastante amplas. Tais equilíbrios podem ser mostrados como equivalentes a alocações de recursos ótimas de Pareto.

Refiro-me a essa teoria como "estática" seguindo Samuelson, apesar da ambiguidade envolvida em especificar uma contrapartida empírica para esse modificador. A ideia subjacente parece ser tirada da física, referindo-se a um sistema "em repouso". Na economia, suponho que tal equilíbrio geral estático corresponda a uma previsão de como uma economia se comportaria caso os choques externos permanecessem fixos por um longo período, de modo que famílias e empresas se ajustassem a enfrentar o mesmo conjunto de preços repetidamente e sintonizassem seu comportamento de acordo. Não é difícil pensar em modificar essa ideia de "repouso" para acomodar mudanças seculares lentas e razoavelmente previsíveis, e essa acomodação será dada como certa no que se segue, como parece ter sido na literatura de economia monetária.

Ora, economias que experimentam ciclos econômicos recorrentes não estão evidentemente "em repouso", de modo que a teoria do equilíbrio geral estático, embora um modelo genuíno no sentido de ser explícito e completo, não é uma boa imitação da realidade para o propósito de entender esses eventos. Para lidar com essa disparidade, o \textit{Foundations} ofereceu uma solução também (embora lá apresentada, parece-me, como uma resposta a um tipo de pergunta inteiramente diferente\footnote{O princípio da correspondência de Samuelson propunha o uso de sua teoria da estabilidade como um critério para ajudar a decidir quais pontos de equilíbrio estacionário poderiam realmente ser observados e quais não: "Quantas vezes o leitor viu um ovo equilibrado em sua extremidade?" Aqui a ideia é claramente decidir quais equilíbrios estáticos de ovos são empiricamente interessantes, e não oferecer um modelo dinâmico empiricamente útil de ovos rolando ou balançando. De fato, a crítica de Gordon e Hynes [10] ao uso da dinâmica de preços de desequilíbrio de Samuelson como uma descrição dos caminhos de preços observados recebeu apoio central de [37].}). Samuelson propôs um modelo dinâmico de ajuste de preços no qual as taxas de variação dos preços oferecidos em cada mercado estavam relacionadas ao nível de "excessos de demanda" em todos os mercados. Qualquer que tenha sido a história ou os objetivos subjacentes deste modelo de dinâmica de preços (e, implicitamente, de dinâmica de quantidades), esta teoria introduziu parâmetros adicionais suficientes (àqueles necessários para descrever gostos e tecnologia) ao sistema de equilíbrio de modo que, dado um choque inicial ao sistema, uma ampla variedade de trajetórias fosse consistente com seu eventual retorno ao equilíbrio.

Esta introdução de parâmetros livres adicionais (àqueles usados para descrever preferências e tecnologia) trazia a promessa de que se poderia construir um sistema teórico cujo ponto estacionário fosse um equilíbrio geral no sentido neoclássico, mas cujos movimentos, fora do equilíbrio, pudessem replicar o comportamento "keynesiano" capturado tão bem pelos modelos econométricos. Explicitar essas conexões tornou-se o principal programa de pesquisa teórica em macroeconomia das décadas de 1940 a 1960. A tarefa parecia envolver de alguma forma motivar a introdução da "moeda" e outros elementos financeiros na teoria do equilíbrio geral "real", modificando o sistema puramente teórico na direção da teoria econométrica aplicada e, simultaneamente, retrabalhando as equações estruturais dos modelos econométricos para esclarecer seus fundamentos teóricos. O objetivo do empreendimento era amplamente aceito como sendo a "unificação" dos dois tipos de teorias nas quais as ideias keynesianas foram traduzidas nas décadas de 1930 e 1940.

Revisar a vasta quantidade de economia útil que resultou dessa tentativa de unificação seria uma tarefa ambiciosa demais para este artigo. Em vez disso, farei algumas observações gerais sobre as atitudes em relação à política de estabilização que foram fomentadas pela síntese neoclássica. Primeiro, como a síntese foi formada pela adição de parâmetros livres a um sistema de equilíbrio geral estático, a classe geral de modelos que ela sugeria admitia uma ampla variedade de possibilidades para o comportamento do ciclo econômico. Assim, parecia um arcabouço suficientemente aberto para conter virtualmente qualquer ponto de vista em relação à política como um caso especial, uma característica atraente para os não dogmáticos. Suspeito que esta seja uma das razões pelas quais economistas como Milton Friedman, que não fizeram uso deste arcabouço, foram tratados com certa impaciência por seus defensores. Por que ele não poderia simplesmente especificar quais valores de parâmetros específicos correspondiam ao caso que ele acreditava ajustar-se aos fatos, deixar que outros fizessem o mesmo e, então, o assunto poderia ser entregue aos econometricistas para uma resolução definitiva?

Segundo, como as flutuações em torno do equilíbrio do sistema representavam um comportamento de desequilíbrio, as proposições de bem-estar padrão poderiam ser aplicadas apenas ao comportamento médio do sistema, e não às flutuações em torno da média. Isso deixava alguém livre para aplicar outros critérios na avaliação de políticas de estabilização: "hiatos" (gaps) em vez de "triângulos", como James Tobin [50] coloca. A ideia geral era usar ferramentas de política para manter a trajetória real do sistema "próxima", em um sentido ou outro, de sua trajetória de equilíbrio. Os proponentes de várias políticas de estabilização estavam, assim, livres do fardo sob o qual o economista do bem-estar comum trabalha — o de "justificar" a intervenção por alguma "falha de mercado" específica e adaptar a natureza da intervenção à natureza da falha. Sob a síntese neoclássica, o ciclo econômico era \textit{definido} como uma falha de mercado, e qualquer política que prometesse mover o sistema em direção ao "equilíbrio de pleno emprego" era vista como uma melhoria.

Tão amplamente endossada foi a ideia geral da síntese neoclássica descrita acima que seus constructos centrais tornaram-se uma taquigrafia comum para descrever "os fatos". Agora, quando o sucesso, real e potencial, desta síntese está novamente em questão, esses constructos não facilitam mais a discussão, mas, pelo contrário, atrapalham. Dois exemplos ilustrarão o que quero dizer.

O primeiro é de James Tobin [50], que, ao resumir o que ele chama de "as proposições centrais da \textit{General Theory}", começa com: "Nas sociedades capitalistas industriais modernas, os preços e salários respondem lentamente ao excesso de demanda ou oferta, especialmente lentamente ao excesso de oferta. Ao longo de um longo curto prazo, os altos e baixos da demanda registram-se no produto; eles estão longe de serem completamente absorvidos pelos preços." Ele prossegue citando evidências para esta proposição a partir do comportamento econômico britânico na década de 1920 e da Grande Depressão nos EUA.

O que entendo que Tobin quer dizer com esta proposição complexa é algo como o seguinte: Se alguém tentasse interpretar a experiência britânica da década de 1920, a americana da década de 1930 e os ciclos econômicos em geral, em termos de um equilíbrio geral estático (permitindo uma tendência secular) e ajuste dinâmico nos preços do tipo descrito por Samuelson, então precisaríamos postular coeficientes "lentos" (digamos, meia-vida de muitos trimestres) descrevendo as respostas de salários e preços ao excesso de oferta. Qualificada desta forma, a afirmação parece-me verdadeira. No entanto, qualificada desta forma, o argumento de Tobin não pode lançar luz sobre a conveniência de tentar explicar os ciclos econômicos dentro do arcabouço fornecido pela síntese neoclássica em oposição ao uso de algum outro arcabouço, que é o assunto pretendido de seu artigo.

Em uma linha semelhante, Franco Modigliani [30] caracteriza o modelo econométrico dos EUA de Thomas Sargent [38] como um no qual os declínios no emprego podem ser explicados apenas por "ataques severos de preguiça contagiosa". Ele prossegue dizendo que "objeções igualmente sérias aplicam-se à modelagem do mercado de commodities feita por Friedman como sendo perfeitamente competitivo... e ao seu tratamento do trabalho como uma commodity homogênea negociada em um mercado de leilão, de modo que, ao salário vigente, nunca há qualquer excesso de demanda pelas empresas ou excesso de oferta pelos trabalhadores".

Como em Tobin em [50], Modigliani está aqui pensando em um equilíbrio competitivo no sentido estático da síntese neoclássica, de modo que uma diminuição de equilíbrio no emprego teria necessariamente que surgir ou de uma tecnologia espontânea ou de uma mudança de gosto ("um ataque de preguiça") e de modo que um modelo no qual os mercados de trabalho são continuamente equilibrados esteja patentemente em contradição com as flutuações observadas no emprego e no desemprego. Mais adiante, ele se refere ao "acordo substancial de que nos Estados Unidos o mecanismo hicksiano [seu termo excessivamente modesto para o que aqui chamo de síntese neoclássica] é bastante eficaz em limitar o efeito de choques e que a resposta de salários e preços ao excesso de demanda e oferta também trabalhará \textit{gradualmente} para eliminar amplamente, se não totalmente, qualquer efeito sobre o emprego". "Essas inferências são apoiadas por simulações com modelos econométricos como o MPS."

Ora, tanto Modigliani quanto Tobin estão, nos artigos dos quais citei, defendendo explicitamente (através da tática consagrada do contra-ataque) seu arcabouço preferido para estudar ciclos econômicos contra alternativas "monetaristas" ou de "expectativas racionais". No entanto, ambos tomam a correção do arcabouço que defendem como dada e o usam como ponto de partida para criticar as alternativas. O fracasso de uma simulação do modelo econométrico de Sargent em reproduzir os resultados de simulações do modelo MPS é visto como evidência de que o modelo de Sargent não "se ajusta aos fatos!" Parece claro que o debate nesse nível não pode fazer as questões avançarem.

\section{RECAPITULAÇÃO E AVALIAÇÃO}

Será útil nesta fase tentar um resumo das teses avançadas nas seções acima. Comecei esboçando uma visão de uma teoria econômica (pela qual entendo uma teoria que se propõe a explicar aspectos específicos observados e ainda não observados do comportamento) como uma economia mecânica análoga. Decorre desta visão que os desenvolvimentos em um campo substantivo particular, como o estudo dos ciclos econômicos, serão fortemente influenciados por desenvolvimentos técnicos em nossa capacidade de construir economias de modelo explícitas, e por desenvolvimentos do mundo real que alteram nossa visão quanto às perguntas que pensamos que tais modelos podem e devem ser capazes de nos ajudar a responder.

Vistos sob essa perspectiva, os principais traços da Revolução Keynesiana e da síntese neoclássica na qual ela evoluiu nos Estados Unidos pareciam ser os seguintes. Eles incluem, em primeira instância, o início da Grande Depressão e a consequente mudança de atenção de explicar um padrão recorrente de altos e baixos para explicar uma economia aparentemente presa em uma baixa interminável. A \textit{General Theory} de Keynes é então vista como, primeiro, um reconhecimento da importância desta mudança de circunstâncias e, segundo, pelo menos como lida por Hicks, Modigliani e outros, como a proposta de uma explicação agregativa simples da determinação do produto e do emprego em um ponto no tempo. Os desenvolvimentos na macroeconomia keynesiana britânica desde a década de 1930 dão, creio eu, uma visão razoavelmente precisa do que teria se tornado a revolução se esses elementos e nenhum outro tivessem sido envolvidos.

Nos Estados Unidos e no continente, dois outros elementos estiveram envolvidos de formas essenciais, ambos de caráter técnico. Um foi o desenvolvimento de descrições estocásticas explícitas de sistemas econômicos. O outro foi o desenvolvimento de uma teoria de equilíbrio geral estático juntamente com uma teoria associada de dinâmica de preços de desequilíbrio. Esses elementos foram rapidamente combinados para fornecer rigor e clareza ao relato de Keynes sobre a determinação do equilíbrio de curto prazo, e para adicionar a essa teoria elementos dinâmicos explícitos, que permitiram que ela se ajustasse às séries temporais reais de uma forma bastante literal. Além disso, eles ofereciam uma promessa definitiva de "unificação" adicional, uma tarefa que ocupou alguns dos melhores talentos da profissão por três décadas.

No final da década de 1960, esse progresso ordenado em direção à unidade foi perturbado por dois desenvolvimentos teóricos. Um deles, o discurso presidencial de Milton Friedman à American Economic Association [8], foi escrito a partir do ponto de vista "monetarista", que continuara a buscar o estudo dos ciclos econômicos ao longo da linha iniciada por Mitchell. O outro, de Edmund Phelps [33] e o subsequente "volume de Phelps" [34], parecia inicialmente uma tentativa de completar a unidade prometida pela síntese neoclássica através da descoberta de uma fundamentação microeconômica para o mercado de trabalho e para o lado da precificação do produto dos modelos padrão. Embora motivados de forma diferente, os artigos de Friedman e Phelps carregavam a clara implicação de que o "excesso de demanda" não era necessário nem suficiente para a inflação de preços ou salários, e que \textit{qualquer} taxa de inflação média era consistente teoricamente com \textit{qualquer} nível de desemprego. Esta conclusão, alcançada através de um raciocínio neoclássico impecável, conflitava com a previsão de um \textit{trade-off} inflação-produto real, que estava no centro de todos os modelos baseados na síntese neoclássica.

Em tentativas de formalizar a hipótese da taxa natural de Friedman-Phelps, logo se descobriu que as formas então convencionais de modelar a formação de expectativas eram centrais para as questões envolvidas e fundamentalmente defeituosas. A hipótese de expectativas racionais de John Muth [31], formulada originalmente para lidar com um conjunto inteiramente diferente de questões substantivas, revelou-se uma forma natural de formalizar os argumentos de Friedman-Phelps. Pesquisas subsequentes em macroeconomia revelaram as implicações abrangentes desta hipótese e a extensão em que ela se mostra subversiva em relação às principais presunções positivas e de política subjacentes à síntese neoclássica.

No momento presente, esses desenvolvimentos retêm novidade suficiente para torná-los difíceis de avaliar com distanciamento. No entanto, acredito que seja possível, e de fato necessário, tentar entendê-los em termos semelhantes aos que usei ao tentar entender desenvolvimentos anteriores na economia monetária e na teoria dos ciclos econômicos. Se foi de fato a conjunção de mudanças no mundo real nas perguntas para as quais as pessoas queriam respostas, juntamente com melhorias técnicas na teoria econômica, que levou à grande reformulação da teoria dos ciclos econômicos que venho chamando aqui de síntese neoclássica, então não é improvável que desenvolvimentos mais recentes possam ser atribuídos de forma semelhante a forças nessas duas categorias.

O evento do mundo real do passado recente que primeiro vem à mente é a combinação de inflação com desemprego acima da média que caracterizou a década de 1970. Embora consistentes com a lógica de Friedman-Phelps, esses eventos foram mal previstos pelos modelos econométricos da safra de 1960. Até que ponto esse erro de previsão deve ser interpretado como um erro "fatal" nos modelos baseados na síntese neoclássica ou simplesmente como um que sugere algumas modificações não é tão fácil de determinar.\footnote{Um argumento em favor da fatalidade está em [26].} A ideia de que virtualmente todo esse período foi caracterizado por "excesso de oferta" e, portanto, que virtualmente toda a inflação deve ser atribuída a "choques de oferta" não parece valer a pena ser levada a sério e ainda não vi um caso quantitativo apresentado para esta posição. (Isso não quer dizer, é claro, que não tenham ocorrido sérios choques de oferta ao longo da década). Por outro lado, alguns modelos recentes que enfatizam o papel de preços nominais fixados contratualmente permitem que os preços respondam à demanda secular ou "antecipada" sem que o "excesso de demanda" necessariamente chegue a emergir, mantendo ao mesmo tempo um papel de "curto prazo" para "excessos de demanda" e oferta.\footnote{Ver [7, 35, 45]. Uma forma de interpretar esses modelos é como tentativas de modificar a dinâmica de preços da síntese neoclássica de modo a permitir que a expansão monetária compartilhe pelo menos parte da culpa pela inflação da década de 1970 com a OPEP, etc. (e pela inflação na Argentina, Chile e inúmeros outros exemplos em que a alta inflação foi associada a produto real e emprego abaixo da tendência). Outra forma é discutida na nota 14.} Talvez estes possam ser interpretados como uma tentativa de reconciliar as experiências da década de 1970 com algumas características fundamentais da síntese neoclássica. Em suma, os eventos da década de 1970 foram provocativos, mas talvez não decisivos.

Uma característica menos espetacular, mas talvez ultimamente mais influente, das séries temporais do pós-Segunda Guerra Mundial tem sido o retorno a um padrão de "ciclos" recorrentes e aproximadamente semelhantes no sentido de Mitchell. Se a magnitude da Grande Depressão desferiu um golpe sério na ideia do ciclo econômico como uma ocorrência repetida do "mesmo" evento, a experiência do pós-guerra restaurou, em algum grau, a respeitabilidade desta ideia. Se a Depressão continua, em alguns aspectos, a desafiar a explicação pela análise econômica existente (como acredito que continue), talvez ela esteja sucumbindo gradualmente à Lei dos Grandes Números.

Essas novas observações têm sido influentes (como novas observações devem ser para pesquisadores empíricos), mas parece-me que as principais influências externas foram, e continuarão a ser, mudanças nos métodos teóricos disponíveis. Na teoria dos ciclos econômicos, parece não ser o problema que muda, mas sim a forma como olhamos para ele. Das mudanças nos métodos, certamente as mais centrais foram os desenvolvimentos do pós-guerra na teoria do equilíbrio geral.

\section{TEORIA DO EQUILÍBRIO GERAL DO PÓS-GUERRA}

A teoria do equilíbrio geral, que Modigliani, Patinkin e outros esperavam integrar com uma teoria operacional do ciclo econômico, não permaneceu congelada na forma que havia assumido nas décadas de 1930 e 1940. De fato, muito do que se desenvolveu desde então foi esboçado em algum detalhe por John Hicks, em \textit{Value and Capital} [15]. Ali, Hicks propôs reinterpretar os problemas de maximização resolvidos por empresas e famílias como envolvendo escolhas sobre sequências de bens datados, com as escolhas de bens futuros específicos interpretadas como planos e com seus preços interpretados como expectativas de preços. Hicks observou que se os bens futuros fossem vistos como contratados antecipadamente, os preços dos bens futuros seriam simplesmente números conhecidos, e o equilíbrio geral em uma economia dinâmica, assim modelada, seria equivalente formalmente ao equilíbrio em um modelo "estático". Hicks acreditava que a presença de incerteza em situações reais tornava um modelo que enfatizasse contratos futuros inaplicável na maioria das situações dinâmicas de interesse real e, por isso, colocou a maior parte de sua ênfase em uma discussão de uma sequência de equilíbrios "à vista" (spot).

Kenneth Arrow [2] e Gerard Debreu [5] observaram que a incerteza poderia ser incorporada na teoria do equilíbrio geral "estática" exatamente pelo mesmo dispositivo que Hicks propôs para incorporar a passagem do tempo: nomeadamente, indexando os bens tanto pela data em que devem ser trocados \textit{quanto} pelo (talvez estocasticamente selecionado) "estado da natureza" contingente ao qual a troca deve ocorrer. Como a de Hicks, a inovação não foi inicialmente uma extensão da teoria do equilíbrio geral no sentido matemático, mas sim a observação de que o \textit{alcance da aplicabilidade} deste corpo de teoria poderia ser vastamente ampliado por alguma engenhosidade na especificação do que se entende por uma mercadoria.

Uma forma de interpretar um equilíbrio de "direitos contingentes" é como uma descrição de uma economia na qual todos os preços contingentes ao estado são determinados antecipadamente, na compensação de um único grande mercado de futuros. Nesta interpretação, os comerciantes individuais podem avaliar as probabilidades da ocorrência de futuros estados da natureza, mas com os preços determinados antecipadamente, a questão das expectativas de preços não surge. Alternativamente, pode-se às vezes (embora certamente nem sempre) pensar em um equilíbrio de direitos contingentes como sendo determinado através de uma sequência de mercados "à vista", nos quais os preços atuais são fixados dadas certas expectativas sobre os preços futuros. Nesta segunda interpretação, precisa-se de um princípio para reconciliar as distribuições de preços implícitas pelo equilíbrio de mercado com as distribuições usadas pelos agentes para formar suas próprias visões do futuro. John Muth [31] observou que o princípio geral da ausência de rendas em equilíbrio competitivo carregava a implicação particular de que essas distribuições não poderiam diferir de forma sistemática. Seu termo para esta última hipótese foi \textit{expectativas racionais}.\footnote{Muth formulou a hipótese de expectativas racionais usando a ideia de equivalência-certeza de Simon [42]-Theil [46], embora sua discussão introdutória deixe claro que a ideia é aplicável em situações onde a equivalência-certeza pode não existir. Sua lógica básica não é difícil de reafirmar dentro do arcabouço de direitos contingentes de Arrow-Debreu: ver [24] para um exemplo.}

Como originalmente proposta por Arrow e Debreu, essa interpretação de direitos contingentes de um modelo de equilíbrio competitivo considerava que toda a informação era simultânea e livremente disponível para todos os negociantes, e muitos resultados importantes (por exemplo, a extensão dos principais teoremas da economia do bem-estar para ambientes incertos) dependem crucialmente desta suposição. Logo foi reconhecido por muitos pesquisadores que a ideia de ver uma mercadoria como uma função de choques determinados estocasticamente é valiosa também em situações nas quais a informação difere de várias maneiras entre os negociantes. De fato, é essa ideia que permite usar a teoria econômica para tornar preciso o que se entende por informação e para determinar como ela é valorizada economicamente.

Quando originalmente proposta, a formulação de direitos contingentes tendia a ser vista como altamente esotérica e remota da prática, sem dúvida porque surgiu na extremidade mais abstrata da disciplina. No entanto, esta formulação absorveu e esclareceu rápida e facilmente uma variedade de resultados especiais na economia da incerteza, facilitando sua unificação e extensão. Ela é agora de uso padrão em virtualmente todos os campos aplicados da economia. Também está em uso, embora eu não diria que ainda seja "padrão", na teoria dos ciclos econômicos.

Seria surpreendente se não fosse assim? A ideia de que elementos especulativos desempenham um papel fundamental nos ciclos econômicos, que esses eventos parecem envolver agentes reagindo a sinais imperfeitos de uma maneira que, após o fato, parece inadequada, tem sido (como observei na seção 3) um lugar-comum na tradição verbal da teoria dos ciclos econômicos pelo menos desde Mitchell [28]. Agora, pela primeira vez, temos à nossa disposição métodos para construir sistemas econômicos de modelo artificial nos quais esses elementos desempenham um papel bem definido. Agora é inteiramente prático ver trajetórias de preços e quantidades que seguem processos estocásticos complicados como "pontos" de equilíbrio em um espaço apropriadamente especificado. Este é um desenvolvimento que fará diferença na forma como pensamos.

Perguntar por que os teóricos monetários da década de 1940 não fizeram uso da visão de equilíbrio de direitos contingentes é, parece-me, como perguntar por que Aníbal não usou tanques contra os romanos em vez de elefantes. Não há razão para ver nossa capacidade de pensar como sendo menos limitada pela tecnologia disponível do que nossa capacidade de agir (se é que essa distinção pode ser defendida). A razão histórica para modelar a dinâmica de preços como respostas a excessos de demanda estáticos não vai além da observação de que os teóricos daquela época não conheciam outra maneira de fazê-lo. É, naturalmente, concebível que teóricos em pleno domínio de novos métodos concluam, no entanto, que existem razões sólidas para continuar a usar a tecnologia mais antiga para alguns propósitos (assim como existem alguns propósitos para os quais continuamos a preferir elefantes a tanques). Isso me parece muito improvável se o propósito for entender os ciclos econômicos. Que um poderoso aparato de construção de modelos especificamente projetado para nos ajudar a lidar com problemas que envolvem escolha sob incerteza deva ser simplesmente ignorado em favor de um aparato mais antigo, que é (na maior parte) incapaz de levar esses problemas em conta,\footnote{Não pretendo sugerir que as considerações de escolha sob incerteza não desempenharam nenhum papel na síntese neoclássica. Tanto a teoria de portfólio quanto a teoria de estoques foram usadas, por exemplo, para motivar hipóteses particulares sobre demandas por moeda e outros ativos [4, 48, 49]. No entanto, o conceito de equilíbrio de mercado usado foi o da teoria do equilíbrio geral estática e determinística.} seria, caso ocorresse, certamente uma refutação decisiva do ponto de vista sobre o desenvolvimento do pensamento econômico que avancei neste artigo.

\section{DESENVOLVIMENTOS FUTUROS}

Mesmo que alguém esteja convencido, como eu, de que os avanços teóricos esboçados na última seção terão um grande impacto em nosso pensamento sobre os ciclos econômicos, é obviamente impossível prever com qualquer precisão a forma que esse impacto tomará. O que já está claro, no entanto, é que certas características de teorias anteriores, adotadas originalmente (conforme argumentei) por conveniência, não são mais analiticamente necessárias. Em particular, é possível construir sistemas em equilíbrio competitivo, em um sentido de direitos contingentes, que exibem uma vasta variedade de comportamentos dinâmicos. A ideia de que um sistema econômico em equilíbrio esteja, de algum modo, "em repouso" é simplesmente um anacronismo.

Para um economista teórico moderno, certamente a maneira mais natural de ler os artigos originais de Friedman e Phelps é como tentativas de conjeturar algumas das propriedades que um modelo de equilíbrio geral bem-sucedido dos ciclos econômicos provavelmente possuiria. De fato, Friedman refere-se à taxa natural de desemprego como "o nível que seria moído pelo sistema walrasiano de equações de equilíbrio geral, desde que nele estivessem embutidas as características estruturais reais dos mercados de trabalho e de commodities", embora ele não seja capaz de colocar tal sistema no papel. Em sua introdução a [34], Phelps esboça um sistema específico de equilíbrio geral com algum detalhe. Grande parte da pesquisa que foi feita desde então pode ser interpretada como esforços para moldar essas ideias mais explicitamente no arcabouço de direitos contingentes.

Nos últimos anos, vários economistas trabalharam para desenvolver o que prefiro chamar de \textit{modelos de equilíbrio de ciclos econômicos}.\footnote{Alguns exemplos iniciais são [20, 21, 41, 38, 3]. Destes, apenas [20] utiliza o formalismo de equilíbrio geral de direitos contingentes em sua totalidade. Os outros usam aproximações lineares, motivadas informalmente por referência aos modelos de direitos contingentes. Parece claro que modelos econometricamente operacionais terão necessariamente que contar com o uso liberal de métodos lineares aproximados.} Estes são modelos que utilizam o ponto de vista de direitos contingentes descrito na última seção de uma forma essencial, e nos quais preços e quantidades são considerados como estando sempre em equilíbrio. Nesses modelos, os conceitos de excessos de demanda e oferta não desempenham papel observacional e não são identificados com nenhuma magnitude observada. Em contraste com os modelos de equilíbrio estático disponíveis na década de 1940, os modelos de equilíbrio desta nova classe parecem se sair tão bem no ajuste de séries temporais quanto os modelos baseados na síntese neoclássica.

Ora, é claro que esses novos modelos de equilíbrio podem, em princípio, ser "sintetizados" com um modelo de ajuste de preços de desequilíbrio do tipo de Samuelson, assim como poderiam ser os modelos de equilíbrio mais antigos.\footnote{Tenho em mente a linha de pesquisa descrita por Malinvaud em [27] e recentemente revisada por Drazen em [6]. A lição de que dois modelos podem estar "próximos" no sentido de se ajustarem aos mesmos dados aproximadamente bem, mas terem implicações radicalmente diferentes para a política, foi trazida de forma contundente em [41] e elaborada de forma mais geral em [39].} Isso deve ser verdade para \textit{qualquer} modelo de equilíbrio. Além disso, como uma síntese desse tipo envolve a adição de parâmetros livres ao modelo, a versão sintetizada não pode se ajustar aos fatos pior do que a versão original de equilíbrio na qual se baseia. Parece-se ser levado, então, não a modelos de equilíbrio como uma classe, mas a uma classe vastamente maior de modelos de desequilíbrio sintetizados. Sinto-me atraído pela visão de que é útil, de uma maneira geral, ser hostil aos teóricos que portam parâmetros livres, de modo que sou simpático à ideia de simplesmente capitalizar essa opinião e chamá-la de um Princípio. Ao avaliar teorias econômicas que pretendem ser úteis na orientação de políticas da forma esboçada em minha seção introdutória, no entanto, existem importantes considerações substantivas que apoiam tal atitude de hostilidade, as quais me parecem colocá-la em uma base mais sólida do que a que pode ser oferecida por um preconceito geral em favor da parcimônia. Tentarei detalhar essas considerações.

Nossa tarefa, tal como a vejo (para reafirmar minha introdução de forma um pouco mais contundente e operacional), é escrever um programa em FORTRAN que aceite regras de política econômica específicas como "entrada" e gere estatísticas de "saída" descrevendo as características operacionais das séries temporais que nos interessam, as quais se prevê que resultem dessas políticas. Por exemplo, gostaríamos de saber qual taxa média de desemprego teria prevalecido desde a Segunda Guerra Mundial nos Estados Unidos se o M1 tivesse crescido a 4 por cento ao ano durante este período, sendo as outras políticas mantidas como foram. (Gostaríamos de saber as respostas para muitas outras perguntas também, mas esta me veio à mente primeiro e dá uma ideia concreta de para onde este argumento se dirige). Deve-se dar por certo, parece claro, que simplesmente tentar várias políticas que possam ser propostas em economias reais e observar o resultado não deve ser levado a sério como um método de solução: Experimentos Sociais em grande escala podem ser instrutivos e admiráveis, mas são melhor admirados à distância. A ideia, se o produto social marginal da economia for positivo, deve ser ganhar alguma confiança de que os componentes do programa são, em algum sentido, confiáveis antes de executá-lo à custa de nossos vizinhos.

Como se ganha confiança desse tipo? Esta é uma questão sobre cuja resposta os economistas estão razoavelmente de acordo, embora eu não me recorde de ter visto a natureza desse acordo articulada. A ideia central é que as respostas \textit{individuais} podem ser documentadas de forma relativamente barata, ocasionalmente por experimentação direta, mas mais comumente por meio do vasto número de instâncias bem documentadas de reações individuais a mudanças ambientais bem especificadas, disponibilizadas "naturalmente" via censos, painéis, outras pesquisas e o método (inapropriadamente difamado como "empirismo casual") de manter os olhos abertos. Sem tais meios de documentar padrões de comportamento, parece claro que o programa em FORTRAN proposto acima não pode ser escrito. Suponha, pelo contrário, que tais meios estejam disponíveis, ou que tenhamos alguma capacidade de prever como o comportamento individual responderá a mudanças especificadas. Como, se é que é possível, tal conhecimento pode ser traduzido em conhecimento sobre a maneira como uma \textit{sociedade} inteira provavelmente reagirá a mudanças em seu ambiente?

Para ser mais concreto, considere a pergunta: Como reagirá um macaco que não foi alimentado por um dia a uma banana lançada em sua jaula? Considero que temos conhecimento previamente estabelecido suficiente sobre o comportamento dos macacos para fazer essa previsão com certa confiança. Agora altere a pergunta para: Como reagirão cinco macacos que não foram alimentados por um dia a uma banana lançada em sua jaula? Esta é uma pergunta inteiramente diferente, na qual o conhecimento das preferências (cada macaco quer o máximo da banana que puder obter) e da tecnologia (o consumo total da banana não pode exceder a unidade) nos dá apenas um começo. Precisamos claramente saber algo sobre a maneira como um grupo de macacos interage, além de suas preferências individuais, para termos qualquer esperança de progresso nesta questão complicada.

Pessoas interessadas na forma como grupos de macacos resolvem problemas de alocação de recursos escassos satisfazem sua curiosidade reunindo grupos de macacos e lançando-lhes recursos escassos. Considerei como dado que nós, economistas, não podemos proceder desta forma, mas a alocação de recursos escassos é algo em que somos admirados por sermos especialistas. A economia é às vezes caracterizada como o trabalho das implicações da ideia de que os indivíduos buscam seu próprio interesse, no entanto, acabamos de observar quão vazia é essa ideia, aplicada isoladamente até ao mais trivial experimento animal. Podemos imaginar que ela ganha poder de alguma forma misteriosa quando aplicada a sociedades humanas com milhões de participantes?

O ingrediente omitido até agora é, naturalmente, a \textit{competição}. Vamos pegar nossa banana, cortá-la em cinco pedaços, dar a cada um dos cinco macacos um pedaço e impor-lhes a regra de que eles podem interagir apenas trocando pedaços de banana por minutos de coçação de costas, a alguma taxa fixa. (Confesso não ter ideia de como tal imposição poderia ser efetuada na prática). Ora, nesta situação, e dadas informações suficientes sobre como os macacos individuais estão dispostos a trocar coçação de costas por consumo de banana, podemos prever o resultado desta interação (preço de equilíbrio e quantidades trocadas), pelo menos dada capacidade computacional suficiente. Observe que, tendo especificado as regras pelas quais a interação ocorre em detalhes, e de uma forma que não introduz \textit{parâmetros livres}, a capacidade de prever o comportamento individual é transformada de forma não experimental na capacidade de prever o comportamento do grupo.\footnote{Este é um caso para o uso da teoria dos jogos explícita em geral, e não para o uso da teoria competitiva em particular. O caso para o uso da teoria competitiva na modelagem de ciclos econômicos seria, se eu fosse desenvolvê-lo aqui, baseado inteiramente na conveniência, ou nos limites impostos a nós pela tecnologia disponível para trabalhar as implicações de outras definições de equilíbrio. Esta ressalva ao texto também deve qualificar a afirmação de que a teoria do equilíbrio exige a introdução de nenhum parâmetro livre além daqueles usados para descrever gostos e tecnologia individuais. Ou seja, na medida em que existe uma latitude considerável quanto a qual conceito de equilíbrio está sendo usado, pode-se pensar em selecionar um deles como a fixação de um "parâmetro livre". Isso não me parece uma questão séria na prática no momento atual, mas pode-se imaginar que venha a ser através de desenvolvimentos futuros nas teorias de jogos não competitivos. Esta observação de que jogos não competitivos podem algum dia provar ser úteis na teoria dos ciclos econômicos, a qual parece difícil de contestar, é às vezes usada para racionalizar modelos totalmente arbitrários não relacionados a qualquer jogo bem definido. Espero que esteja claro que esta nota não se destina a ser uma defesa desta prática.}

Enfatizei que é a hipótese do equilíbrio competitivo que permite que o comportamento do grupo seja previsto a partir do conhecimento das preferências individuais e da tecnologia sem a adição de quaisquer parâmetros livres. Isso precisa de mais ilustração, em contextos mais próximos de nossos interesses do que os experimentos com animais. O emprego e os salários nominais são, em um sentido imediato, determinados por algumas interações muito complicadas no mercado de trabalho envolvendo empregados e empregadores. É possível, sabemos, imitar o resultado agregado desta interação razoavelmente bem de uma forma de equilíbrio competitivo, na qual os salários e as horas de trabalho são gerados pela interação de famílias e empresas "representativas".\footnote{Ver [25], embora partes deste estudo precisem de atualização. Uma replicação moderna deste estudo utilizaria ideias em [40, 12].} Os parâmetros neste modelo descrevem ou a disposição das famílias em substituir bens e lazer contemporânea e intertemporalmente ou a tecnologia disponível para as empresas.

Também é possível, à maneira da síntese neoclássica, ajustar um modelo bastante diferente — uma curva de Phillips — a esses mesmos resultados agregados. Aqui também se utiliza uma descrição paramétrica (diferente, é claro, neste caso) das preferências e da tecnologia e, \textit{adicionalmente}, um parâmetro descrevendo a velocidade com que um "leiloeiro" ajusta o salário nominal aos excessos de oferta e demanda. Ora, a introdução de um leiloeiro fictício não é um defeito desta segunda maneira de ver as coisas, em relação à primeira. Todos os modelos são ficções de acordo com o ponto de vista que estou adotando e, em qualquer caso, um leiloeiro é pressuposto no primeiro modelo também, mas um operando tão rapidamente que não é notado. Nem pode ser uma desvantagem desta segunda forma de modelar a determinação de salários e emprego que ela não possa se ajustar aos dados tão bem quanto a primeira: a adição de um parâmetro livre não pode prejudicar nesse sentido.

A desvantagem do segundo modelo é esta: não há como obter informações sobre a taxa na qual o salário é ajustado, esse parâmetro adicionado, exceto observando o sistema inteiro em operação. Se este parâmetro mudar em reação a mudanças em outras partes do sistema (como sabemos que de fato acontece), não há como prever a natureza dessas respostas sem experimentar com o sistema como um todo. No entanto, é precisamente a tentativa de evitar ter que fazer isso que nos leva a usar a teoria econômica em primeiro lugar.

No caso da explicação de equilíbrio da determinação de salários e emprego, parâmetros que descrevem o grau de substituibilidade intertemporal fazem o trabalho (no sentido empírico) do parâmetro que descreve o comportamento do leiloeiro no modelo da curva de Phillips. Sobre esses parâmetros, temos uma riqueza de dados disponíveis a baixo custo, provenientes de informações de coortes de censo, de dados de painel que descrevem as reações de famílias individuais a uma variedade de mudanças nas condições de mercado, e assim por diante. Em princípio (e talvez em breve, na prática, pois há muita pesquisa promissora ocorrendo justamente sobre este tópico\footnote{Por exemplo, [9, 13, 19, 1]. Apresso-me a acrescentar que esses estudos e outros relacionados não foram motivados ou planejados como apoio ou confirmação dos resultados de Rapping e os meus, nem está claro neste estágio que estejam levando a estimativas consistentes com as nossas. Meu ponto é simplesmente que os modelos agregados de equilíbrio carregam implicações para uma variedade de outros tipos de dados, levantando a possibilidade de confirmação independente (e contradição) de estimativas a partir de séries temporais agregativas.}), esses parâmetros cruciais podem ser estimados de forma independente a partir de dados individuais e agregados. Se assim for, saberemos o que os parâmetros agregados significam, nós os entenderemos em um sentido em que os parâmetros de ajuste de desequilíbrio \textit{nunca} serão entendidos. É precisamente por isso que nos importamos com as "fundamentações microeconômicas" das teorias agregadas.\footnote{Os modelos citados na nota 6 constituem uma terceira classe, intermediária àquelas acabadas de discutir? Não creio que sim. Se a duração do contrato nesses artigos for vista simplesmente como um parâmetro livre (como a nota 6 os interpreta), então este parâmetro é tão ininteligível (no sentido deste parágrafo) quanto aquele que descreve a rapidez do ajuste de um leiloeiro. Se, por outro lado, a duração do contrato for vista como emergindo de um problema de decisão resolvido pelos agentes, então esses modelos, assim elaborados, seriam modelos de equilíbrio (com um espaço de mercadorias diferente daqueles discutidos acima) e não serviriam necessariamente para reforçar o ponto de vista em relação à política adotado pela síntese neoclássica. Se esta observação deve ser tomada como uma crítica severa aos modelos baseados em contratos depende, naturalmente, da visão de cada um quanto à probabilidade de sermos capazes de explicar os ciclos econômicos usando nenhum parâmetro de ajuste livre. Certamente esta questão deve ser considerada como aberta no presente, e independentemente de como venha a ser resolvida, deve ser certamente o caso de que modelos com um ou dois parâmetros de ajuste livres representem progresso analítico sobre modelos com dezenas ou centenas deles.}

Pesquisadores familiarizados com o trabalho atual aludido no parágrafo anterior apreciarão a extensão em que ele descreve esperanças para o futuro, e não realizações passadas. Essas esperanças poderiam, sem esforço, ser descritas como esperanças por um tipo de unificação, não muito diferente em espírito da esperança de unificação que informou a síntese neoclássica. O que tentei fazer acima foi enfatizar o caráter empírico (em oposição ao estético) dessas esperanças, para tentar entender como tal evidência quantitativa sobre o comportamento, que podemos razoavelmente esperar obter na sociedade tal como ela existe agora, poderia concebivelmente ser transformada em informação quantitativa sobre o comportamento de sociedades imaginadas, diferentes em aspectos importantes de qualquer uma que já tenha existido. Isso pode parecer uma forma intimidante e ambiciosa de declarar o objetivo de um subcampo aplicado de uma ciência marginalmente respeitável, mas existe uma maneira menos ambiciosa de descrever o objetivo da teoria dos ciclos econômicos?

\section{OBSERVAÇÕES FINAIS}

Este artigo foi uma tentativa de entender e esclarecer a natureza e as origens de alguns desenvolvimentos recentes na teoria dos ciclos econômicos. Ao adotar um ponto de vista histórico, pode ser que eu tenha inadvertidamente adotado a tática marxista de descrever o que eu gostaria que acontecesse como algo que a História já ordenou. Eu, certamente, enfatizei o que me parece ser a importância decisiva das melhorias no equipamento analítico que temos à nossa disposição, mas com a intenção de sublinhar a expansão de nossas oportunidades que essas melhorias oferecem, não as direções particulares que elas ditam. Por esta razão, tentei evitar reivindicar demais para os exemplos particulares de modelos de equilíbrio que existem agora. Não há sentido em deixar que os primeiros passos, embora tentativos e, espero, promissores, endureçam em posições que devem ser defendidas a todo custo.

Se uma abordagem histórica não pode garantir a capacidade de prever o futuro, ela me parece ajudar a distinguir aqueles elementos no pensamento passado que permanecem úteis daqueles que não permanecem. A síntese neoclássica surgiu, como toda economia útil, de um compromisso entre o que gostaríamos de ter sabido e o que os métodos à nossa disposição pareciam tornar possível saber. Nada poderia ser mais prejudicial ao uso produtivo de métodos desenvolvidos mais recentemente do que ver as categorias e constructos que foram produzidos por este compromisso como restrições à forma como pensamos sobre os ciclos econômicos hoje.

\section*{REFERÊNCIAS CITADAS}

\begin{enumerate}
    \item Altonji, Joseph G., e Orley C. Ashenfelter. "Wage Movements and the Labor Market Equilibrium Hypothesis." Princeton University, Industrial Relations Section, Working Paper No. 130, Novembro 1979.
    \item Arrow, Kenneth J. "The Role of Securities in the Optimal Allocation of Risk-Bearing." \textit{Review of Economic Studies}, 31 (Abril 1964), 91--96.
    \item Barro, Robert J. "Rational Expectations and the Role of Monetary Policy." \textit{Journal of Monetary Economics}, 2 (Janeiro 1976), 1--32.
    \item Baumol, William. "The Transactions Demand for Money --- An Inventory Theoretic Approach." \textit{Quarterly Journal of Economics}, 66 (Novembro 1952), 545--56.
    \item Debreu, Gerard. \textit{Theory of Value}. New Haven, Conn.: Yale University Press, 1959.
    \item Drazen, Allan. "Recent Developments in Macroeconomic Disequilibrium Theory." \textit{Econometrica}, 48 (Março 1980), 283--306.
    \item Fischer, Stanley. "Long-Term Contracts, Rational Expectations, and the Optimal Money Supply Rule." \textit{Journal of Political Economy}, 85 (Fevereiro 1977), 191--206.
    \item Friedman, Milton. "The Role of Monetary Policy." \textit{American Economic Review}, 58 (Março 1968), 1--17.
    \item Ghez, Gilbert R., e Gary S. Becker. \textit{The Allocation of Time and Goods over the Life Cycle}. New York: National Bureau of Economic Research, 1975.
    \item Gordon, Donald F., e J. Allan Hynes. "On the Theory of Price Dynamics." In \textit{Macroeconomic Foundations of Employment and Inflation Theory}, editado por Edmund S. Phelps, pp. 369--93, New York: Norton, 1970.
    \item Haberler, Gottfried. \textit{Prosperity and Depression}. Geneva: League of Nations, 1937.
    \item Hansen, Lars P., e Thomas J. Sargent. "Formulating and Estimating Dynamic Linear Rational Expectations Models." \textit{Journal of Economic Dynamics and Control}, 2 (1980), 7--46.
    \item Heckman, James J. "Longitudinal Studies in Labor Economics: A Methodological Review." Working Paper, University of Chicago, Setembro 1978.
    \item Hicks, John R. "Mr. Keynes and the 'Classics': A Suggested Interpretation." \textit{Econometrica}, 5 (Abril 1937), 147--59.
    \item \rule{1cm}{0.4pt}. \textit{Value and Capital: An Inquiry Into Some Fundamental Principles of Economic Theory}. Oxford: Clarendon Press, 1939.
    \item Keynes, John M. \textit{A Treatise on Money}. New York: Harcourt Brace and Co., 1930.
    \item \rule{1cm}{0.4pt}. \textit{The General Theory of Employment, Interest and Money}. London: Macmillan, 1936.
    \item Leijonhufvud, Axel. \textit{On Keynesian Economics and the Economics of Keynes}. New York: Oxford, 1968.
    \item Lillard, Lee A., e Robert J. Willis. "Dynamic Aspects of Earning Mobility." \textit{Econometrica} (Setembro 1978), 985--1012.
    \item Lucas, Robert E., Jr. "Expectations and the Neutrality of Money." \textit{Journal of Economic Theory}, 4 (Abril 1972), 103--24.
    \item \rule{1cm}{0.4pt}. "An Equilibrium Model of the Business Cycle." \textit{Journal of Political Economy}, 83 (Dezembro 1975), 1113--44.
    \item \rule{1cm}{0.4pt}. "Econometric Policy Evaluation: A Critique." \textit{Journal of Monetary Economics}, 2, Supplement (1976), Carnegie-Rochester Conference Series, Vol. 1.
    \item \rule{1cm}{0.4pt}. "Understanding Business Cycles." \textit{Journal of Monetary Economics}, Supplement (1977), Carnegie-Rochester Conference Series, Vol. 5.
    \item Lucas, Robert E. Jr., e Edward C. Prescott. "Investment Under Uncertainty." \textit{Econometrica}, 39 (Setembro 1971), 659--81.
    \item Lucas, Robert E. Jr., e Leonard A. Rapping. "Real Wages, Employment, and the Price Level." \textit{Journal of Political Economy}, 77 (Outubro 1969), 721--54.
    \item Lucas, Robert E. Jr., e Thomas J. Sargent. "After Keynesian Macroeconomics." In \textit{After the Phillips Curve: Persistence of High Inflation and High Unemployment}, Conference Series no. 19, pp. 49--72. Boston, Mass.: Federal Reserve Bank of Boston.
    \item Malinvaud, Edmund. \textit{The Theory of Unemployment Reconsidered}. New York: Wiley, 1977.
    \item Mitchell, Wesley C. \textit{Business Cycles}. Berkeley, Calif.: University of California Press, 1913.
    \item Modigliani, Franco. "Liquidity Preference and the Theory of Interest and Money." \textit{Econometrica}, 12 (Janeiro 1944), 45--88.
    \item \rule{1cm}{0.4pt}. "The Monetarist Controversy or, Should We Forsake Stabilization Policies?" \textit{American Economic Review}, 67 (Março 1977), 1--19.
    \item Muth, John F. "Rational Expectations and the Theory of Price Movements." \textit{Econometrica}, 29 (Julho 1961), 315--35.
    \item Patinkin, Don. \textit{Money, Interest, and Prices}. Second edition. New York: Harper and Row, 1965.
    \item Phelps, Edmund S. "Money Wage Dynamics and Labor Market Equilibrium." \textit{Journal of Political Economy}, 76 (Julho/Agosto 1968), 687--711.
    \item Phelps, Edmund S., et al. \textit{Microeconomic Foundations of Employment and Inflation Theory}. New York: Norton, 1970.
    \item Phelps, Edmund S., e John B. Taylor. "Stabilizing Powers of Monetary Policy under Rational Expectations." \textit{Journal of Political Economy}, 85 (Fevereiro 1977), 163--89.
    \item Samuelson, Paul A. \textit{Foundations of Economic Analysis}. Cambridge, Mass.: Harvard University Press, 1947.
    \item \rule{1cm}{0.4pt}. "Proof that Properly Anticipated Prices Fluctuate Randomly." \textit{Industrial Management Review}, 6 (1965), 41--49.
    \item Sargent, Thomas J. "A Classical Macroeconomic Model for the United States." \textit{Journal of Political Economy}, 84 (Abril 1976), 207--38.
    \item \rule{1cm}{0.4pt}. "The Observational Equivalence of Natural and Unnatural Rate Theories of Macroeconomics." \textit{Journal of Political Economy}, 84 (Junho 1976), 631--40.
    \item \rule{1cm}{0.4pt}. "Estimation of Dynamic Labor Demand Schedules under Rational Expectations." \textit{Journal of Political Economy}, 86 (Dezembro 1978), 1009--44.
    \item Sargent, Thomas J., e Neil Wallace. "'Rational' Expectations, The Optimal Monetary Instrument, and the Optimal Money Supply Rule." \textit{Journal of Political Economy}, 83 (Abril 1975), 241--54.
    \item Simon, Herbert A. "Dynamic Programming under Uncertainty with a Quadratic Criterion Function." \textit{Econometrica}, 24 (Janeiro 1956), 74--81.
    \item \rule{1cm}{0.4pt}. \textit{The Sciences of the Artificial}. Cambridge, Mass.: MIT Press, 1969.
    \item Slutsky, Eugenio. "The Summation of Random Causes as the Source of Cyclic Processes." \textit{Econometrica}, 5 (Abril 1937), 105--46.
    \item Taylor, John G. "Estimation and Control of a Macroeconomic Model with Rational Expectations." \textit{Econometrica}, 47 (Setembro 1979), 1267--86.
    \item Theil, Henri. "A Note on Certainty Equivalence in Dynamic Programming." \textit{Econometrica}, 25 (1957), 346--49.
    \item Tinbergen, Jan. \textit{Business Cycles in the United States of America 1919-1932. Statistical Testing of Business Cycle Theories}, Vol. 2. Geneva, League of Nations, 1939.
    \item Tobin, James. "The Interest Elasticity of the Transactions Demand for Cash." \textit{Review of Economics and Statistics}, 38 (Agosto 1956), 241--47.
    \item \rule{1cm}{0.4pt}. "Liquidity Preference as Behavior towards Risk." \textit{Review of Economic Studies}, 25 (Fevereiro 1958), 65--86.
    \item \rule{1cm}{0.4pt}. "How Dead is Keynes?" \textit{Economic Inquiry}, 16 (1977), 459--68.
\end{enumerate}

\end{document}